% ============================================================================
% COURS 04 — EXERCICES DE CINÉMATIQUE — PARTIE B
% Source : 5-Cinematique2022.docx
% ============================================================================

\subsection{Manipulateur de stèles de marbre}

\textit{Source : Jacques Le Goff.}

Le manipulateur permet l'élévation, l'inclinaison et l'orientation d'une
stèle pouvant atteindre $400\ \mathrm{kg}$. On étudie une combinaison
des mouvements d'élévation et d'inclinaison.

\begin{center}
\TLCineImage[.31\linewidth]{images/image197.pdf}\hfill
\TLCineImage[.31\linewidth]{images/image198.png}\hfill
\TLCineImage[.31\linewidth]{images/image200.png}
\end{center}

\begin{enumerate}
  \item Construire la chaîne de solides et les figures de changement de base.
  \item Déterminer, par dérivation du vecteur position, la vitesse du point $O_8$ par rapport à $R_0$.
  \item Retrouver le résultat par composition des vitesses.
  \item Déterminer l'accélération du point $O_8$.
  \item Identifier les termes dus à l'élévation, à l'inclinaison et à leur couplage.
\end{enumerate}

\subsection{Robot de manutention à parallélogramme}

Le robot déplace son poignet dans un plan à l'aide de deux moteurs. Les
solides 1, 2, 3 et 4 forment une structure en parallélogramme
déformable.

\begin{center}
\TLCineImage[.46\linewidth]{images/image222.png}\hfill
\TLCineImage[.46\linewidth]{images/image223.jpeg}
\end{center}

\begin{enumerate}
  \item Établir les relations entre les bases imposées par le parallélogramme.
  \item Déterminer les torseurs cinématiques des pivots commandés.
  \item Composer les torseurs pour obtenir les mouvements 3/0 et 4/0.
  \item Déterminer la vitesse d'un point du poignet par le champ des vitesses.
  \item Décrire sa trajectoire lorsque chacun des moteurs est successivement bloqué.
\end{enumerate}

\subsection{Monte-charge}

Une charge est entraînée par un câble et une poulie. Un tour moteur
correspond à une montée de $10\ \mathrm{cm}$. Le cycle comporte une
accélération jusqu'à un capteur situé à $0{,}5\ \mathrm m$, un palier à
$0{,}2\ \mathrm{m\,s^{-1}}$ jusqu'à $1{,}5\ \mathrm m$, puis un
freinage. La montée dure $15\ \mathrm s$.

\TLCineImage[.48\linewidth]{images/image238.png}

\begin{enumerate}
  \item Déterminer l'accélération de la première phase et les équations horaires.
  \item Calculer sa durée.
  \item Déterminer les équations horaires et la durée de la phase uniforme.
  \item Déterminer la décélération de la dernière phase.
  \item Calculer la distance totale de montée et tracer la vitesse moteur.
\end{enumerate}

\subsection{Robot de soudage quatre axes}

Le robot réalise un cordon linéaire puis un dégagement rapide de
l'organe terminal. Chaque axe possède son propre actionneur.

\begin{center}
\TLCineImage[.31\linewidth]{images/image239.jpeg}\hfill
\TLCineImage[.31\linewidth]{images/image241.png}\hfill
\TLCineImage[.31\linewidth]{images/image242.png}
\end{center}

\begin{enumerate}
  \item Déterminer les vitesses des points caractéristiques des solides 1 et 2.
  \item Donner le torseur du solide 2 par rapport au bâti, réduit en $O_2$.
  \item Déterminer la vitesse de l'extrémité de soudage.
  \item Établir les relations entre les paramètres pour obtenir une trajectoire rectiligne imposée.
  \item Lorsque deux axes sont bloqués, déterminer l'accélération pendant le dégagement rapide.
\end{enumerate}

\subsection{Scénic contre Ferrari de Formule 1}

Les deux véhicules effectuent : accélération de 0 à
$100\ \mathrm{km\,h^{-1}}$, maintien pendant $10\ \mathrm s$, puis
freinage jusqu'à l'arrêt. Le Scénic atteint $100\ \mathrm{km\,h^{-1}}$
en $13{,}4\ \mathrm s$ et freine sur $46\ \mathrm m$ ; la F1 atteint
la même vitesse en $2{,}3\ \mathrm s$ et freine sur $14\ \mathrm m$.

\begin{center}
\TLCineImage[.45\linewidth]{images/image254.jpeg}\hfill
\TLCineImage[.45\linewidth]{images/image255.jpeg}
\end{center}

\begin{enumerate}
  \item Établir $a(t)$, $v(t)$ et $x(t)$ pour les trois phases de chaque véhicule.
  \item Calculer la distance totale parcourue.
  \item Calculer le temps et la distance nécessaires pour atteindre la vitesse maximale de chaque véhicule.
  \item Comparer les accélérations et décélérations moyennes.
\end{enumerate}

\subsection{Manipulateur de collecteur d'échappement}

\textit{Source : Mines PTSI 2006.}

Le manipulateur assiste un opérateur pour déplacer un collecteur de
$17\ \mathrm{kg}$. Un vérin équilibre le poids ; les autres mouvements
sont guidés manuellement.

\begin{center}
\TLCineImage[.34\linewidth]{images/image256.jpeg}\hfill
\TLCineImage[.28\linewidth]{images/image257.png}\hfill
\TLCineImage[.30\linewidth]{images/image258.png}
\end{center}

\begin{enumerate}
  \item Compléter le schéma cinématique et le paramétrage.
  \item Déterminer les vecteurs rotation associés aux différentes liaisons.
  \item Écrire les torseurs cinématiques au point de réduction adapté.
  \item Déterminer la vitesse d'un point $P$ du collecteur par composition.
  \item Discuter les positions singulières et les directions de vitesse accessibles.
\end{enumerate}

\subsection{Robot de pick and place}

Le schéma cinématique d'un robot de prise-dépose comporte plusieurs
pivots et une glissière.

\begin{center}
\TLCineImage[.43\linewidth]{images/image264.png}\hfill
\TLCineImage[.43\linewidth]{images/image265.png}
\end{center}

\begin{enumerate}
  \item Déterminer les vecteurs rotation des solides successifs.
  \item Écrire les torseurs cinématiques des liaisons.
  \item Déterminer les vitesses de points par déplacement des torseurs.
  \item Construire les figures de changement de base.
  \item Déterminer la vitesse et l'accélération de l'effecteur.
  \item Identifier les mouvements de translation possibles entre deux solides.
\end{enumerate}

\subsection{Shimmy d'une roue de train d'atterrissage}

L'avion 1 est en translation rectiligne uniforme par rapport au sol 0.
L'étrier 2 pivote autour d'un axe vertical et porte la roue 3, de rayon
$R$, roulant sur la piste au point $A$. Le mouvement de shimmy est une
oscillation de l'étrier autour de son pivot.

\TLCineImage[.66\linewidth]{images/image287.png}

\begin{enumerate}
  \item Écrire la condition de roulement sans glissement en $A$.
  \item Appliquer la composition des vitesses pour déterminer la vitesse du centre $O_3$.
  \item Relier la vitesse du centre et celle du point de contact par le champ des vitesses.
  \item Déduire la relation entre la vitesse de l'avion, la rotation de la roue et l'oscillation de l'étrier.
\end{enumerate}

\subsection{Canne robotisée}

Une roue motorisée de rayon $R$ déplace une canne inclinée d'un angle
$\theta(t)$. La longueur $AH=l(t)$ varie afin de maintenir la poignée
$H$ à la hauteur constante $h_0$.

\begin{center}
\TLCineImage[.43\linewidth]{images/image293.png}\hfill
\TLCineImage[.43\linewidth]{images/image294.png}
\end{center}

\begin{enumerate}
  \item Écrire la fermeture géométrique $\overrightarrow{IA}+\overrightarrow{AH}+\overrightarrow{HI}=\vec0$.
  \item Projeter pour relier $l(t)$, $\theta(t)$, $R$ et $h_0$.
  \item Écrire la condition de roulement sans glissement en $I$.
  \item Déterminer $\vec V(H\in3/2)$ et les vecteurs rotation.
  \item Composer les vitesses pour obtenir $\vec V(H\in3/0)$.
  \item Déduire la vitesse de la roue assurant la vitesse de déplacement imposée sans mouvement parasite de la main.
\end{enumerate}

\subsection{Dépose de bagage automatisée}

Le système déplace et bascule un bagage entre différents convoyeurs. La
cinématique doit garantir la continuité des vitesses et limiter les
accélérations transmises au bagage.

\begin{center}
\TLCineImage[.29\linewidth]{images/image295.png}\hfill
\TLCineImage[.29\linewidth]{images/image296.png}\hfill
\TLCineImage[.34\linewidth]{images/image297.png}
\end{center}

\begin{enumerate}
  \item Identifier les mouvements élémentaires et les trajectoires des points du bagage.
  \item Écrire les torseurs des liaisons puis les composer.
  \item Déterminer la vitesse et l'accélération d'un point caractéristique.
  \item Vérifier les contraintes de temps de transfert et de maintien du bagage.
\end{enumerate}

\subsection{Véhicule autoguidé hospitalier}

Le VAG doit s'arrêter depuis $v_{\max}=1{,}2\ \mathrm{m\,s^{-1}}$ sur
une distance maximale $d_f=0{,}2\ \mathrm m$. La vitesse décroît
linéairement pendant le freinage.

\begin{center}
\TLCineImage[.40\linewidth]{images/image298.pdf}\hfill
\TLCineImage[.48\linewidth]{images/image299.png}
\end{center}

\begin{enumerate}
  \item Exprimer la durée du freinage en fonction de $d_f$ et $v_{\max}$.
  \item Calculer la décélération constante nécessaire.
  \item Tracer les courbes $v(t)$, $a(t)$ et $x(t)$.
  \item Vérifier la compatibilité avec la condition de non-dérapage.
\end{enumerate}

\subsection{Synthèse : choix d'une méthode}

Pour chacun des cas suivants, indiquer la méthode la plus efficace avant
de réaliser le calcul demandé :
\begin{enumerate}
  \item vitesse de deux points d'un même solide ;
  \item vitesse d'un point dans une chaîne de trois solides ;
  \item dérivée d'un vecteur unitaire d'une base mobile ;
  \item relation cinématique imposée par un contact sans glissement ;
  \item déplacement associé à une loi de vitesse en trapèze ;
  \item accélération d'un point déjà connu par son vecteur vitesse.
\end{enumerate}
